% ================================================================
%  Annexures — KANDA Project
% ================================================================
\chapter*{ANNEXURE}
\markboth{Annexure}{}
\addcontentsline{toc}{chapter}{Annexure}
\label{ch:annexures}
\setcounter{section}{0}
\renewcommand{\thesection}{\arabic{section}}

%-------------------------------------------------------------------
\section{Pin Configuration and Wiring}
\label{sec:pinout}

\textbf{Table~\ref{tab:esp32_pins}} documents the ESP32 GPIO assignments used in the firmware.

\begin{table}[H]
\centering
\caption[ESP32 GPIO pin assignments]{\textbf{ESP32 GPIO pin assignments}}
\label{tab:esp32_pins}
\begin{tabularx}{\textwidth}{|l|l|X|}
\hline
\textbf{GPIO} & \textbf{Function} & \textbf{Notes} \\
\hline
5 & Front ultrasonic TRIG & 3.3V output \\
\hline
18 & Front ultrasonic ECHO & 5V$\to$3.3V divider \\
\hline
19 & Left ultrasonic TRIG & 3.3V output \\
\hline
21 & Left ultrasonic ECHO & 5V$\to$3.3V divider \\
\hline
22 & Right ultrasonic TRIG & 3.3V output \\
\hline
23 & Right ultrasonic ECHO & 5V$\to$3.3V divider \\
\hline
25 & Motor A IN1 & TB6612FNG direction \\
\hline
26 & Motor A IN2 & TB6612FNG direction \\
\hline
27 & Motor A PWM & LEDC channel 0, 1 kHz \\
\hline
32 & Motor B IN1 & TB6612FNG direction \\
\hline
33 & Motor B IN2 & TB6612FNG direction \\
\hline
14 & Motor B PWM & LEDC channel 1, 1 kHz \\
\hline
SDA (21) & OLED SDA & I\textsuperscript{2}C, address 0x3C \\
\hline
SCL (22) & OLED SCL & Shared with Right TRIG (multiplexed) \\
\hline
TX (1) & UART TX to Pi & 115200 baud \\
\hline
RX (3) & UART RX from Pi & 115200 baud \\
\hline
\end{tabularx}
\end{table}

%-------------------------------------------------------------------
\section{Firmware Source Code (Excerpt)}
\label{sec:fw_code}

The following excerpt shows the core control loop of the ESP32 firmware including sensor reading, mode switching, and command execution.

\begin{lstlisting}[language=C++, caption={ESP32 firmware main loop (simplified)}, label={lst:fw_main}]
void loop() {
  // 1. Read ultrasonic sensors
  float front = readUltrasonic(TRIG_F, ECHO_F);
  float left  = readUltrasonic(TRIG_L, ECHO_L);
  float right = readUltrasonic(TRIG_R, ECHO_R);

  // 2. Send telemetry to Pi
  Serial.printf("F:%.1f L:%.1f R:%.1f -> %s\n",
                front, left, right, modeLabel);

  // 3. Check for AI commands
  if (Serial.available()) {
    String line = Serial.readStringUntil('\n');
    if (executeAiCommand(line)) {
      currentMode = AI_MODE;
      lastAiTime = millis();
    }
  }

  // 4. Mode timeout check
  if (currentMode == AI_MODE &&
      millis() - lastAiTime > AI_TIMEOUT_MS) {
    currentMode = AUTO_MODE;
  }

  // 5. Execute based on mode
  if (currentMode == AUTO_MODE) {
    autoObstacleAvoidance(front, left, right);
  }

  // 6. Safety reflex (overrides everything)
  if (front < EMERGENCY_STOP_CM) {
    stopMotors();
  }

  updateOLED(currentMode, front);
  delay(LOOP_DELAY_MS);
}
\end{lstlisting}

%-------------------------------------------------------------------
\section{Intelligence Layer Source Code (Excerpt)}
\label{sec:ai_code}

The following excerpt shows the Python orchestrator control cycle on the Raspberry Pi.

\begin{lstlisting}[language=Python, caption={Intelligence layer main loop (simplified)}, label={lst:ai_main}]
def run_cycle(bridge, agent, validator):
    telemetry = bridge.read_telemetry()
    if telemetry is None:
        bridge.send_command("stop", 0)
        return

    prompt = build_prompt(telemetry)
    response = agent.get_action(prompt)
    safe_cmd = validator.validate(response)
    bridge.send_command(
        safe_cmd["action"],
        safe_cmd["speed"]
    )

def main():
    bridge = UARTBridge(config.SERIAL_PORT)
    bridge.connect()
    agent = TaskAgent()
    validator = SafetyValidator()

    signal.signal(signal.SIGINT,
                  lambda *_: shutdown(bridge))

    while running:
        try:
            run_cycle(bridge, agent, validator)
        except Exception as e:
            logger.error(f"Cycle error: {e}")
            bridge.send_command("stop", 0)
        time.sleep(config.LOOP_INTERVAL_SEC)
\end{lstlisting}

%-------------------------------------------------------------------
\section{Raspberry Pi to ESP32 Wiring Diagram}
\label{sec:wiring}

\begin{table}[H]
\centering
\caption[UART wiring between Pi and ESP32]{\textbf{UART wiring between Raspberry Pi and ESP32}}
\label{tab:uart_wiring}
\begin{tabularx}{\textwidth}{|l|l|X|}
\hline
\textbf{Raspberry Pi} & \textbf{ESP32} & \textbf{Notes} \\
\hline
GPIO 14 (TX) & GPIO 3 (RX) & Pi TX to ESP32 RX \\
\hline
GPIO 15 (RX) & GPIO 1 (TX) & ESP32 TX to Pi RX \\
\hline
GND & GND & Common ground required \\
\hline
\end{tabularx}
\end{table}

\begin{table}[H]
\centering
\caption[Power distribution]{\textbf{Power distribution}}
\label{tab:power}
\begin{tabularx}{\textwidth}{|l|l|X|}
\hline
\textbf{Component} & \textbf{Supply} & \textbf{Source} \\
\hline
Raspberry Pi 4 & 5V 3A & Buck converter from LiPo \\
\hline
ESP32 DevKit & 5V (USB or Vin) & Buck converter from LiPo \\
\hline
TB6612FNG motors & 7.4V direct & LiPo battery (2S) \\
\hline
HC-SR04 sensors & 5V & Shared 5V rail \\
\hline
SSD1306 OLED & 3.3V & ESP32 3.3V regulator \\
\hline
\end{tabularx}
\end{table}
